% Part: model-theory
% Chapter: models-of-arithmetic
% Section: non-standard-models

\documentclass[../../../include/open-logic-section]{subfiles}

\begin{document}

\olfileid{mod}{mar}{nst}
\section{అప్రామాణిక నమూనాలు}

\begin{explain}
$\Lang{L_A}$కు గల ఒక \tetoken{నిర్మాణం}{!!a{structure}},
$\Struct{N}$తో సమరూపమైతే దాన్ని ప్రామాణికం అంటాం. ఒక
\tetoken{నిర్మాణం}{!!a{structure}}~$\Struct{N}$తో సమరూపం కాకపోతే
దాన్ని అప్రామాణికం అంటారు.
\end{explain}

\begin{defn}
$\Lang{L_A}$కు గల \tetoken{నిర్మాణం}{!!{structure}}~$\Struct{M}$,
$\Struct{N}$తో సమరూపం కాకపోతే \emph{అప్రామాణికం}. ఏదో
$n\in\Nat$కు $\Value{\num{n}}{M}$తో సమానమైన
$x\in\Domain{M}$ \tetoken{మూలకాలను}{!!{element}s}
($\Struct{M}$ యొక్క) \emph{ప్రామాణిక సంఖ్యలు} అంటారు; అలా కానివాటిని
\emph{అప్రామాణిక సంఖ్యలు} అంటారు.
\end{defn}

\begin{explain}
\olref[stm]{prop:standard-domain} వల్ల, $\Lang{L_A}$కు గల ప్రతి
ప్రామాణిక \tetoken{నిర్మాణం}{!!{structure}} ప్రామాణిక
\tetoken{మూలకాలను}{!!{element}s} మాత్రమే కలిగి ఉంటుంది. అందువల్ల ఒక
అప్రామాణిక \tetoken{నిర్మాణం}{!!{structure}} కనీసం ఒక అప్రామాణిక
మూలకాన్ని కలిగి ఉండాలి. నిజానికి, ఒక అప్రామాణిక
\tetoken{మూలకం}{!!{element}} ఉనికే ఆ
\tetoken{నిర్మాణం}{!!{structure}} అప్రామాణికమని నిర్ధారిస్తుంది.
\end{explain}

\begin{prop}
$\Lang{L_A}$కు గల \tetoken{నిర్మాణం}{!!a{structure}}~$\Struct{M}$
ఒక అప్రామాణిక సంఖ్యను కలిగి ఉంటే, $\Struct{M}$ అప్రామాణికం.
\end{prop}

\begin{proof}
అలా కాదనుకుందాం; అంటే $\Struct{M}$ ప్రామాణికమే అయినా అప్రామాణిక
సంఖ్య $x$ను కలిగి ఉందనుకుందాం. $g\colon\Nat\to\Domain{M}$ ఒక
సమరూపతగా ఉండనివ్వండి. $n$పై ఆగమనం ద్వారా
$g(\Value{\num{n}}{N})=\Value{\num{n}}{M}$ అని సులభంగా తెలుస్తుంది.
అంటే $g$, $\Struct{N}$లోని ప్రామాణిక సంఖ్యలను $\Struct{M}$లోని
ప్రామాణిక సంఖ్యలకు పంపుతుంది. $\Struct{M}$ ఒక అప్రామాణిక సంఖ్యను
కలిగి ఉంటే $g$ \tetoken{సంగ్రస్త}{!!{surjective}} కాలేదు; ఇది
పరికల్పనకు విరుద్ధం.
\end{proof}

\begin{prob}
$\Th{Q}$లో ఈ స్వీకృతాలు ఉన్నాయని గుర్తుచేసుకోండి:
\begin{align*}
& \lforall[x][\lforall[y][(\eq[x'][y'] \lif \eq[x][y])]] \tag{$!Q_1$}\\
& \lforall[x][\eq/[\Obj 0][x']] \tag{$!Q_2$}\\
& \lforall[x][(\eq[x][\Obj 0] \lor \lexists[y][\eq[x][y']])] \tag{$!Q_3$}
\end{align*}
కింది షరతులు తీరేలా
\tetoken{నిర్మాణాలు}{!!{structure}s}~$\Struct{M_1}$,
$\Struct{M_2}$, $\Struct{M_3}$ ఇవ్వండి:
\begin{enumerate}
\item $\Sat{M_1}{!Q_1}$, $\Sat{M_1}{!Q_2}$, $\Sat/{M_1}{!Q_3}$;
\item $\Sat{M_2}{!Q_1}$, $\Sat/{M_2}{!Q_2}$, $\Sat{M_2}{!Q_3}$; మరియు
\item $\Sat/{M_3}{!Q_1}$, $\Sat{M_3}{!Q_2}$, $\Sat{M_3}{!Q_3}$.
\end{enumerate}
ప్రతి దానికీ $\Assign{\Obj{0}}{M_i}$, $\Assign{\prime}{M_i}$లను
మాత్రమే నిర్దేశించాల్సి ఉంటుందని స్పష్టమే.
\end{prob}

\begin{explain}
$\Lang{L_A}$కు అప్రామాణిక
\tetoken{నిర్మాణాలను}{!!{structure}s} ఇవ్వడం చాలా సులభం. ఉదాహరణకు,
\tetoken{వ్యక్తి క్షేత్రం}{!!{domain}}~$\Int$ గల నిర్మాణాన్ని తీసుకొని
తార్కికేతర సంకేతాలన్నింటికీ మామూలు అర్థనిర్దేశాలు ఇవ్వండి. ఋణ సంఖ్యలు
ఏ $n$కూ $\num{n}$ విలువలు కావు కాబట్టి ఈ నిర్మాణం అప్రామాణికం.
అయితే $\Struct{N}$ సత్యం చేసే వాక్యాలనే ఇది కూడా సత్యం చేయాలనే
అర్థంలో, ఇది అంకగణితపు \emph{నమూనా} కాదు. ఉదాహరణకు,
$\lforall[x][\eq/[x'][\Obj{0}]]$ అసత్యం. అయినప్పటికీ, సంహతత్వ
సిద్ధాంతాన్ని ఉపయోగించి అంకగణితానికి అప్రామాణిక నమూనాలు ఉన్నాయని
సులభంగా నిరూపించవచ్చు.
\end{explain}

\begin{prop}
$\Th{TA}=\Setabs{!A}{\Sat{N}{!A}}$ అనేది $\Struct{N}$ యొక్క
సిద్ధాంతం అనుకుందాం. $\Th{TA}$కు
\tetoken{ఒక లెక్కించదగిన}{!!a{enumerable}} అప్రామాణిక నమూనా ఉంది.
\end{prop}

\begin{proof}
$\Lang{L_A}$ను కొత్త \tetoken{స్థిరసంకేతం}{!!{constant}}~$c$తో
విస్తరించి, ఈ \tetoken{వాక్యాల}{!!{sentence}s} సమితిని పరిశీలించండి:
\[
\Gamma = \Th{TA} \cup \{\eq/[c][\num{0}], \eq/[c][\num{1}],
\eq/[c][\num{2}], \dots\}
\]
$\Gamma$ యొక్క ఏ నమూనా~$\Struct{M^c}$ అయినా, ప్రతి $n\in\Nat$కూ
$x\neq\Value{\num{n}}{M}$ అయ్యే అప్రామాణిక
\tetoken{మూలకం}{!!a{element}}~
$x=\Assign{c}{M^c}$ను కలిగి ఉంటుంది. అలాగే
$\Th{TA}\subseteq\Gamma$ కాబట్టి $\Sat{M^c}{\Th{TA}}$ అని స్పష్టమే.
\sourcecorrection{OLTEMODARI-005}{కొత్త స్థిరసంకేతం $c$కు అర్థనిర్దేశం
విస్తృత నిర్మాణం $\Struct{M^c}$లోనే ఉంటుంది; ఆంగ్ల మూలం పొరపాటుగా
$\Assign{c}{M}$ అంటుంది. $c$ను మరిచిన తరువాతి సంకుచిత రూపంలో అది
నిర్వచితం కాదు కాబట్టి నిర్మాణ సూచికను సరిచేశాం.}
$c$ను మరిచిపోయి $\Struct{M^c}$ను $\Lang{L_A}$కు గల
\tetoken{నిర్మాణం}{!!a{structure}}~$\Struct{M}$గా మార్చినా, దాని
వ్యక్తి క్షేత్రం అదే అప్రామాణిక $x$ను కలిగి ఉంటుంది; అలాగే
$\Sat{M}{\Th{TA}}$. రెండవ విషయం $\Th{TA}$లో $c$ రాదు కాబట్టి
నిర్ధారితమవుతుంది. కనుక $\Gamma$కు ఒక నమూనా ఉందని చూపితే చాలు.

$\Gamma$కు ఒక నమూనా ఉందని చూపడానికి సంహతత్వ సిద్ధాంతాన్ని వాడతాం.
$\Gamma$ యొక్క ప్రతి పరిమిత ఉపసమితి సంతృప్తిపరచదగినదైతే
$\Gamma$ కూడా అలాగే ఉంటుంది. ఏదైనా పరిమిత
$\Gamma_0\subseteq\Gamma$ను పరిశీలించండి. $\Gamma_0$లో
$\Th{TA}$కు చెందిన కొన్ని \tetoken{వాక్యాలు}{!!{sentence}s},
$\eq/[c][\num{n}]$ రూపంలోని కొన్ని వాక్యాలు ఉంటాయి; రెండవ రకమైనవి
పరిమిత సంఖ్యలో మాత్రమే. అవేవీ లేకపోతే $\Struct{N}$ను
$\Assign{c}{N_0}=0$గా విస్తరించిన $\Struct{N_0}$,
$\Gamma_0$ను సంతృప్తిపరుస్తుంది. అవి ఉంటే,
$\eq/[c][\num{k}]\in\Gamma_0$ అయ్యే అతిపెద్ద సంఖ్య $k$ అనుకుందాం.
$\Assign{c}{N_k}=k+1$గా $\Struct{N}$ను విస్తరించి
$\Struct{N_k}$ను నిర్వచించండి. అప్పుడు $\Sat{N_k}{\Gamma_0}$:
$!A\in\Th{TA}$ అయితే $\Struct{N_k}$, $c$ విషయంలో తప్ప మిగతా
అన్ని విధాలా $\Struct{N}$లాగే ఉంటుంది, $!A$లో $c$ రాదు కాబట్టి
$\Sat{N_k}{!A}$. అలాగే $n\le k$, $\Value{c}{N_k}=k+1$ కాబట్టి
$\Sat{N_k}{\eq/[c][\num{n}]}$. ఈ విధంగా $\Gamma$ యొక్క ప్రతి
పరిమిత ఉపసమితి సంతృప్తిపరచదగినది.
\sourcecorrection{OLTEMODARI-006}{పరిమిత ఉపసమితిలో
$\eq/[c][\num{n}]$ రూపపు వాక్యమేదీ లేకపోతే “అతిపెద్ద $k$”
ఉండదు. ఆంగ్ల మూల నిరూపణ వదిలేసిన ఈ సందర్భానికి
$\Assign{c}{N_0}=0$ అనే ఏ విస్తరణైనా సరిపోతుందని చేర్చాం.}
కాబట్టి సంహతత్వం $\Gamma$కు ఒక నమూనాను ఇస్తుంది. విస్తరించిన భాష
లెక్కించదగినది కాబట్టి, అధోముఖ లొవెన్‌హైమ్--స్కోలెమ్ సిద్ధాంతం ద్వారా
ఆ నమూనాను లెక్కించదగినదిగా ఎంచుకోవచ్చు.
\sourcecorrection{OLTEMODARI-007}{ప్రతిపాదన లెక్కించదగిన నమూనాను
కోరుతుంది; ఆంగ్ల మూల నిరూపణ సంహతత్వం ద్వారా ఏదో నమూనా ఉందని మాత్రమే
చూపి ఆ షరతును నిర్ధారించదు. అదే ఘనీభూత మూలంలోని సమాంతర వివరణ
లెక్కించదగిన భాషకు అధోముఖ లొవెన్‌హైమ్--స్కోలెమ్ సిద్ధాంతాన్ని
స్పష్టంగా ఉపయోగిస్తుంది; ఆ తప్పిపోయిన దశను చేర్చాం.}
\end{proof}

\end{document}
